\section{Fourier表示与离散变换}
\label{sec:fourier-discrete}

到目前为止出现了两种频率表示：周期函数的Fourier系数是整数下标的序列，全空间函数的Fourier变换是连续频率的函数。实际计算中还有第三种：有限个数组成的向量，用有限个频率表示。三者都用复指数作坐标，但输入对象、频率集合和成立的意义都不同。本节先把三种对象并列，再说明从一种转到另一种时，需要哪些额外的约定，以及每种约定改变了什么。

\paragraph{离散Fourier变换}

有限向量的情形最简单。设$N\geq1$，把$\bm x\in\mathbb C^N$的分量记为$x[0],\ldots,x[N-1]$，并把下标看作模$N$的整数，即首尾相接。此时允许的频率只有$N$个，对应的复指数为$e^{2\pi\mathrm ikn/N}$，$k=0,\ldots,N-1$。

\begin{definition}[离散Fourier变换]
\label{def:fourier-dft}
$\bm x\in\mathbb C^N$的\term{离散Fourier变换}（discrete Fourier transform, DFT）为
\[
\widehat x[k]=\sum_{n=0}^{N-1}x[n]e^{-2\pi\mathrm ikn/N},\qquad k=0,\ldots,N-1 .
\]
\end{definition}

逆变换与能量关系都来自单位根的正交性：对整数$m$，$\sum_{n=0}^{N-1}e^{2\pi\mathrm imn/N}$在$N\mid m$时等于$N$，否则是公比$q=e^{2\pi\mathrm im/N}\neq1$且$q^N=1$的几何级数，等于$(1-q^N)/(1-q)=0$。由此
\[
x[n]=\frac1N\sum_{k=0}^{N-1}\widehat x[k]e^{2\pi\mathrm ikn/N},\qquad
\sum_{n=0}^{N-1}\lvert x[n]\rvert^2=\frac1N\sum_{k=0}^{N-1}\lvert\widehat x[k]\rvert^2 .
\]
这里所有的和都是有限和，不存在收敛问题，变换是$\mathbb C^N$上的可逆线性映射，$N^{-1/2}$倍的变换矩阵是酉矩阵。常数$1/N$的位置同样是约定，与第~\ref{sec:fourier-l2}节的表格同理。

\paragraph{三种表示的对象}

下表把三种表示并列。它们的差别首先在输入：全空间上的函数、一个周期上的函数、$N$个数。输入不同，频率集合与成立的意义也随之不同。

\begin{table}[!htb]
  \centering
  \caption{三种Fourier表示的输入对象、频率集合与成立意义。表中周期情形写成一般周期$T$，第~\ref{sec:fourier-series}节取$T=2\pi$。}
  \label{tab:fourier-three-representations}
  \small
  \begin{tabularx}{\linewidth}{@{}l>{\raggedright\arraybackslash}X>{\raggedright\arraybackslash}X>{\raggedright\arraybackslash}X@{}}
    \toprule
    表示 & 输入对象 & 频率集合 & 成立的意义\\
    \midrule
    Fourier变换 & $\mathbb R^d$上的$L^1$或$L^2$函数 & $\xi\in\mathbb R^d$ & $L^1$：处处有值；$L^2$：等价类\\
    Fourier级数 & 周期$T$的$L^1$或$L^2$函数 & $k/T$，$k\in\mathbb Z$ & $L^2$展开；逐点另需条件\\
    离散变换 & $\mathbb C^N$中的向量 & $k$，$k=0,\ldots,N-1$（模$N$） & 有限和，精确可逆\\
    \bottomrule
  \end{tabularx}
\end{table}

这三种对象不能直接互相代替。一个全空间函数没有周期，不能直接取Fourier系数；一个周期函数有无穷多个系数，不能直接作$N$点变换；$N$个数不是函数，没有连续频率。从一种表示转到另一种，必须先把对象本身转过去，下面三个操作各完成一种转换。

\paragraph{周期化：从全空间到周期}

把全空间函数变成周期函数的自然做法是把所有平移叠加起来。对$f\in L^1(\mathbb R)$，定义它的\term{周期化}（periodization）
\[
(Pf)(x)=\sum_{n\in\mathbb Z}f(x+n).
\]
由Tonelli定理，$\int_0^1\sum_n\lvert f(x+n)\rvert\,\mathrm dx=\int_{\mathbb R}\lvert f\rvert<\infty$，所以级数对几乎每个$x$绝对收敛，$Pf$以$1$为周期且在一个周期上可积。

\begin{theorem}[周期化的系数]
\label{thm:fourier-periodization}
若$f\in L^1(\mathbb R)$，则$Pf$按周期$1$计算的Fourier系数为
\[
\int_0^1(Pf)(x)e^{-2\pi\mathrm ikx}\,\mathrm dx=\widehat f(k),\qquad k\in\mathbb Z .
\]
\end{theorem}

\begin{proof}
上面的绝对收敛估计允许交换求和与积分。$e^{-2\pi\mathrm ikx}$以$1$为周期，所以
\begin{equation*}
\begin{aligned}
\int_0^1\sum_nf(x+n)e^{-2\pi\mathrm ikx}\,\mathrm dx&=\sum_n\int_n^{n+1}f(y)e^{-2\pi\mathrm iky}\,\mathrm dy\\
&=\int_{\mathbb R}f(y)e^{-2\pi\mathrm iky}\,\mathrm dy .
\end{aligned}
\end{equation*}
\end{proof}

定理说明，周期化后的Fourier系数恰是原变换在整数频率上的取值，连续频率被取样成了离散频率。代价在原函数一侧：$Pf$把相距整数的各段叠在一起。若$f$支撑在长度不超过$1$的区间内，各段不重叠，$Pf$在一个周期上就是$f$，没有信息损失；否则不同位置的值混在一起，由$Pf$不能还原$f$。以$f=e^{-\lvert x\rvert}$为例，$Pf$的系数为$2/(1+4\pi^2k^2)$，而$Pf$本身在每一点都叠加了无穷多个平移。由这一系数等式还可以得到逐点的Poisson求和公式$\sum_nf(n)=\sum_k\widehat f(k)$，但那需要额外的正则性，这里不展开。

\paragraph{取样：从周期函数到有限向量}

把周期函数变成$N$个数的自然做法是在一个周期内等距取样。设$g$以$1$为周期，Fourier系数$\widehat g(m)=\int_0^1g(x)e^{-2\pi\mathrm imx}\,\mathrm dx$绝对可和；由定理~\ref{thm:fourier-uniform}（换成周期$1$），二次连续可微就足够。此时Fourier级数一致收敛到$g$，在每一点成立。取样$x[n]=g(n/N)$，$n=0,\ldots,N-1$。

\begin{theorem}[取样的离散变换]
\label{thm:fourier-sampling-aliasing}
在上述条件下，
\[
\widehat x[k]=N\sum_{j\in\mathbb Z}\widehat g(k+jN),\qquad k=0,\ldots,N-1 .
\]
\end{theorem}

\begin{proof}
$x[n]=\sum_m\widehat g(m)e^{2\pi\mathrm imn/N}$，级数绝对收敛，可以与关于$n$的有限和交换：
\[
\widehat x[k]=\sum_m\widehat g(m)\sum_{n=0}^{N-1}e^{2\pi\mathrm i(m-k)n/N}=N\sum_{m\equiv k\ (\mathrm{mod}\ N)}\widehat g(m),
\]
最后一步用了单位根的正交性。
\end{proof}

于是$\widehat x[k]/N$不是单个系数$\widehat g(k)$，而是所有与$k$模$N$同余的频率的系数之和：取样后，频率$k$与$k\pm N,k\pm2N,\ldots$无法区分。这一现象称为\term{混叠}（aliasing）。例如$g(x)=\cos(6\pi x)$只有$\widehat g(\pm3)=1/2$，取$N=4$时，$m=3$落在$k=3$，$m=-3\equiv1$落在$k=1$，样本的离散变换在$k=1,3$处等于$2$，与$\cos(2\pi x)$的样本完全相同；事实上$\cos(3\pi n/2)=\cos(\pi n/2)$对每个整数$n$成立。反过来，若$g$是三角多项式，且非零系数的频率$m$都满足$\lvert m\rvert<N/2$，则每个同余类中至多有一个非零系数，离散变换精确给出全部系数。

\paragraph{截断：从全空间到有限记录}

实际只能观察全空间函数的一段。把$f$截断到长度$L$的区间上，相当于乘以方框$B_L=\mathbf 1_{[-L/2,L/2]}$。乘积在频域中对应卷积：对$f,g\in L^2$，把第~\ref{sec:fourier-convolution}节$L^2$卷积的对应关系用于$\widehat f,\widehat g$，再用$\mathcal F\mathcal F=R$，可得$\widehat{fg}=\widehat f*\widehat g$。所以截断后的变换是$\widehat f$与$\widehat{B_L}(\xi)=\sin(\pi L\xi)/(\pi\xi)$的卷积，原来集中在某个频率的成分会向邻近频率扩散。

由此可以看清把连续变换交给有限计算的完整链条：先截断（频域中与$\widehat{B_L}$卷积），再以$L$为周期周期化（频率被取样到$k/L$），最后在一个周期内取$N$个样本（频率模$N/L$混叠）。每一步都引入一种约定，也带来一种相应的误差，离散变换给出的数只在这些约定下近似连续变换。这些误差何时可以忽略、怎样控制，取决于原函数的衰减与光滑性，后者正是第~\ref{sec:fourier-operations}节的微分规则所描述的。

本节只核对三种表示的对象和它们之间的精确关系。把这些表示用于具体信号，比较信号的频谱，属于\BookChapterRef{chap:signal-analysis}{卷九“信号的表示与比较”章}；卷积作为系统对信号的作用以及滤波，属于\BookChapterRef{chap:signal-systems}{卷九“信号系统”章}；由样本恢复信号的条件与方法，属于\BookChapterRef{chap:sampling-signal-recovery}{卷九“采样与信号恢复”章}。
